% !TEX root = ../main.tex
\section{\sys}
\label{sec:design}

Section~\ref{sec:mot} leads to three requirements. First, capacity must follow each chunk's load without control flow:
the program must choose at run time which expert a lane serves, and bind capacities to a quantity that is stable across
chunks. Second, aggregation must stay efficient whatever shapes the capacities take: keep each device's expert rows on that
device, sum each token's rows on one core, and cost in proportion to the routed rows, not to lanes$\,\times\,T$. Third, both
must be static graphs of operations that ahead-of-time targets provide: reductions, a sort, gathers with run-time indices
and one cross-device sum. This section states the method independent of the accelerator; \S\ref{sec:impl} summarizes how
the Cloud AI 100 compiler shaped its realization.

\begin{figure*}[t]
\centering
% Figure: one MoE layer under RankTier. Bars: sorted token counts of SoC 2 in layer 20 of the held-out 512-token chunk.
\begin{tikzpicture}[x=1in, y=1in, font=\footnotesize, >={Stealth[length=3.6pt, width=3pt]},
  box/.style={draw=black!70, rounded corners=2pt, align=center, inner sep=3pt, fill=white, line width=0.5pt},
  flow/.style={->, line width=0.7pt, draw=black!75},
  num/.style={circle, fill=black!80, text=white, inner sep=0.7pt, font=\scriptsize\bfseries},
  head/.style={anchor=west, font=\footnotesize\bfseries, inner sep=1pt}]
\definecolor{tierA}{RGB}{31,90,166}\definecolor{tierB}{RGB}{79,143,209}\definecolor{tierC}{RGB}{169,200,234}
% frame for the per-SoC steps 1-3
\draw[rounded corners=4pt, fill=black!4, draw=black!35, line width=0.5pt] (0.74,0.02) rectangle (5.12,2.02);
\node[anchor=south east, font=\scriptsize\itshape, text=black!55, inner sep=1pt] at (5.10,0.04) {on each SoC};
% router
\node[box] (router) at (0.30,1.00) {router\\[-1pt]top-8\\[-1pt]of 128};
\draw[flow] (router.east) -- (0.86,1.00);
% ---- (1) rank the SoC's 32 experts
\node[num] at (0.88,1.86) {1}; \node[head] at (0.96,1.86) {rank by load};
\foreach [count=\i] \c in {176,115,80,62,56,53,52,49,39,37,32,30,26,24,15,14,13,10,9,9,8,2,2,1,1,0,0,0,0,0,0,0} {
  \pgfmathsetmacro{\x}{0.90 + (\i-1)*0.031}
  \pgfmathsetmacro{\h}{max(0.006, \c/176*1.05)}
  \ifnum\i<9 \def\col{tierA}\else\ifnum\i<17 \def\col{tierB}\else\def\col{tierC}\fi\fi
  \fill[\col] (\x,0.40) rectangle ++(0.025,\h);
}
\draw[black!60, line width=0.4pt] (0.89,0.40) -- (1.90,0.40);
\node[anchor=north west, font=\scriptsize, text=black!65, inner sep=1pt] at (0.88,0.38) {rank 1};
\node[anchor=north east, font=\scriptsize, text=black!65, inner sep=1pt] at (1.91,0.38) {32};
\node[anchor=north, font=\scriptsize, align=center, text=black!75, inner sep=1pt] at (1.40,0.24) {tokens per expert, sorted};
% ---- (2) tiers; weights gathered by rank
\node[num] at (2.10,1.86) {2}; \node[head] at (2.18,1.86) {fixed-capacity tiers};
\node[box, minimum width=1.30in, fill=black!2, font=\scriptsize] (bank) at (2.85,1.56) {weights of all 128 experts (replicated)};
\draw[flow] (1.60,0.62) to[out=80, in=180] node[pos=0.55, left, font=\scriptsize, inner sep=2pt] {$\pi_c$} (bank.west);
\draw[flow] (2.85,1.46) -- node[right, font=\scriptsize, inner sep=2pt] {gather by rank} (2.85,1.32);
\def\L{0.03125}
\fill[tierA] (2.20,1.05) rectangle (3.50,1.30);
\fill[tierB] (2.20,0.80) rectangle (2.525,1.05);
\fill[tierC] (2.20,0.30) rectangle (2.3625,0.80);
\foreach \k in {1,...,31} {\draw[white, line width=0.25pt] (2.20,{0.30+\k*\L}) -- ({\k>23 ? 3.50 : (\k>15 ? 2.525 : 2.3625)},{0.30+\k*\L});}
\draw[dashed, black!55, line width=0.5pt] (2.20,0.30) rectangle (3.50,1.30);
\node[fill=tierA, text=white, font=\scriptsize, inner sep=1pt] at (2.85,1.175) {tier 1: 8 lanes $\times\,T$};
\node[anchor=west, font=\scriptsize, inner sep=1pt] at (2.55,0.925) {tier 2: 8 lanes $\times\,C_2$};
\node[anchor=west, font=\scriptsize, inner sep=1pt] at (2.39,0.68) {tier 3: 16 lanes $\times\,C_3$};
\node[anchor=south east, font=\scriptsize\itshape, text=black!55, align=right, inner sep=1.5pt] at (3.50,0.31) {capacity-$T$ padding:\\[-2pt]32 lanes $\times\,T$};
\node[anchor=north, font=\scriptsize, text=black!75, inner sep=1pt] at (2.85,0.24) {rows per lane (capacity)};
% ---- (3) token-owned combine
\node[num] at (3.72,1.86) {3}; \node[head] at (3.80,1.86) {token-owned combine};
\node[box, text width=1.18in, font=\scriptsize] (comb) at (4.38,0.95) {each core takes $T/16$ tokens, gathers the $\le 8$ rows of each token that this SoC computed, and sums them\\[2pt] partial $P_c\in\mathbb{R}^{T\times d}$};
\draw[flow] (3.51,0.95) -- (comb.west);
% ---- (4) sum over SoCs
\node[num] at (5.28,1.86) {4}; \node[head] at (5.36,1.86) {sum over SoCs};
\node[box, text width=1.12in, font=\scriptsize] (sum) at (6.02,1.25) {{\footnotesize $(P_0{+}P_1){+}(P_2{+}P_3)$}\\[2pt] elementwise, lowered as a reduce-scatter};
\draw[flow] (comb.east) -- (5.30,0.95) |- (sum.west);
\draw[flow] (sum.south) -- (6.02,0.62) node[below, font=\scriptsize, inner sep=2pt] {next layer};
\end{tikzpicture}

\vskip -0.05in
\caption{One MoE layer under \sys. (1) Each SoC counts the tokens routed to its 32 experts and sorts them (bars: SoC~2 in
layer 20 of the held-out 512-token chunk; colors mark the tiers). (2) Three tiers of lanes with compile-time capacities
($C_2=128$, $C_3=64$ at $T=512$) serve the experts by rank; their weights are gathered from a bank of all 128 experts
through the rank permutation $\pi_c$. The filled area is the padded work, against the dashed capacity-$T$ layout.
(3) Each core gathers and sums the rows of its own tokens. (4) One elementwise sum over SoCs completes the layer.}
\label{fig:overview}
\vskip -0.1in
\end{figure*}

\subsection{A Static Model of an MoE Layer}
\label{sec:model}

A device $c$ owns a fixed set $\mathcal{E}_c$ of experts, 32 of the 128 on each SoC here. A static program computes them in
\emph{lanes}: a lane is one expert's three GEMMs over $C$ token rows gathered through a token table, and lanes of equal
capacity that run together form a stage. The boundary between compile time and run time is
sharp. Fixed at compile time: the number of tiers, their lanes and capacities, the shapes of every tensor, the weight bank
and the expert-to-device placement. Decided by data at run time: which expert each lane serves, which tokens fill its
rows, which weights it gathers, and which rows each token's aggregation reads. A chunk routes $n_e$
of its $kT$ assignments to expert $e$ and $N_c=\sum_{e\in\mathcal{E}_c}n_e$ to device $c$, and a lane that serves $e$ drops
tokens unless $n_e\le C$. A lane's cost has a weight part, its expert's matrices streamed from memory, and a row part
that grows with the rows it computes: all $C$ of them without control flow, or $R\lceil n_e/R\rceil$ on a target that
predicates row tiles of $R$ rows (\S\ref{sec:skip}). Either way the capacity, fixed before the routing is known, bounds a
lane's work.

\subsection{Rank-Bound Capacities}
\label{sec:sort}\label{sec:tiers}

After the router, each device counts its experts' tokens and sorts the counts, which gives a permutation $\pi_c$ from rank
to expert with $n_{\pi_c(1)}\ge n_{\pi_c(2)}\ge\dots$. The device's lanes form tiers: tier $i$ is a stage of $L_i$ lanes
at capacity $C_i$ that covers ranks $a_i,\dots,a_i+L_i-1$, with $\sum_iL_i=|\mathcal{E}_c|$, and lane $j$ of tier $i$
serves expert $\pi_c(a_i+j)$. The lane gathers that expert's weights through $\pi_c$ and its token table from tables
that the routing computes for all experts in their fixed order (Figure~\ref{fig:overview}). Every expert keeps exactly one
lane, so a tiering changes only the padding: the same experts compute the same tokens, and as long as every expert fits
its lane, the output equals that of any other tiering, capacity $T$ everywhere included (bit-identical in our
implementation). A device provisions $\sum_iL_iC_i$ rows per layer against $|\mathcal{E}_c|\,T$.

\textbf{Why ranks.} The order statistics of per-device load are bounded for a structural reason. The $r$ most loaded
experts of device $c$ each hold at least $n_{\pi_c(r)}$ assignments, and together at most $N_c$, so
\begin{equation}
n_{\pi_c(r)}\;\le\;N_c/r,\qquad N_c\;\le\;\min(k,|\mathcal{E}_c|)\,T .
\label{eq:bound}
\end{equation}
$N_c$ sums over all of the device's experts, $kT|\mathcal{E}_c|/E=2T$ on average here and at most 1.9$\times$ that in our
256- and 512-token captures, so every rank beyond the first few is bounded well below $T$ whatever the workload. A capacity
attached to an expert must instead cover the worst case of that one expert, which approaches $T$ because routing is
skewed and which experts are hot changes with the workload (\S\ref{sec:mot-pad}). Eq.~(\ref{eq:bound}) holds for any
routing; that the per-rank maxima are also stable across workloads, while expert identities are not, is a property of the
model, which we measured on held-out documents of eight workloads (Figure~\ref{fig:rank}, \S\ref{sec:eval-gen}).

\subsection{Calibrating the Tiers}
\label{sec:calib}

\textbf{Capacities.} The first tier has capacity $T$, since rank 1 can receive every token. For $i>1$, $C_i$ is the largest
load observed at rank $a_i$, the tier's first rank, over calibration chunks, times a headroom $h$, rounded up to the
target's row granularity $g$ and capped at $T$; since ranks are sorted, rank $a_i$ bounds the whole tier. We calibrate
with $h=1.28$ and $g=16$ on routing captures of the eight workloads: the 128-token prompts of \S\ref{sec:mot-pad}, 56
chunks of 256 tokens and 20 of 512.

\textbf{Structure.} The number and widths of the tiers trade padded rows against per-stage cost. Rows alone favor many
tiers: under our calibration at $T=256$ and 512, three tiers need 65--67\% of the rows of a two-tier split at $T$ and $T/8$, four 52--58\% and
five 48--56\%. But each tier is a stage with a cost of its own, a weight pipeline and one aggregation gather, and a target
may restrict the widths a stage can take. We fix the structure by measurement, not by a model: on Cloud AI 100, where
widths must divide 16 and stages run one after another (\S\ref{sec:impl}), we compared two, three and four tiers and
deploy three tiers of 8, 8 and 16 lanes per SoC at every chunk length, with capacities $T$, $C_2$ and $C_3$ of 128/48/16,
256/64/32 and 512/128/64 for $T=128$, 256 and 512: 12--13$T$ rows per SoC against $18T$ for two
tiers and $32T$ for capacity-$T$ padding. This structure is not optimal everywhere: at $T=256$ two tiers are 2.5\% faster
in the full model, although three save a third of the rows (\S\ref{sec:eval-pad}). Choosing the structure per chunk
length from measured stage costs remains open.

\subsection{Overflow Detection and Recovery}
\label{sec:overflow}

The lower tiers are calibrated, not dropless, so a deployment must detect overflow and recover from it exactly. Detection
is free: the program already computes every lane's token count, and the host compares it with the lane's capacity after
each chunk. Recovery reruns the chunk on a capacity-$T$ variant at the same positions. Prefill writes the KV cache only at
the chunk's own positions, so the rerun overwrites every entry the failed attempt wrote, later chunks see only the rerun's
cache, and the result is that of an overflow-free run. The variant can be a second profile of the same program, since
tier capacities read from the lengths of small row inputs let the compiler specialize one program per profile over the
same weights and cache, or a second program kept resident, which must first rebuild the cache from the earlier chunks.
With overflow probability $p$ a chunk costs $t_{\mathrm{tier}}+p\,t_T$ on
average, so the tiered program's own latency matters most (\S\ref{sec:eval-chunks}).

\subsection{Ownership-Preserving Token Aggregation}
\label{sec:red}

After the expert GEMMs, each token's $k$ weighted expert rows must be summed. Device $c$ computes the experts
$\mathcal{E}_c$, and the result is
\begin{equation}
P_c[t]=\sum_{j:\,e_{t,j}\in\mathcal{E}_c} O_c[\mathrm{row}(t,j)],\qquad y[t]=\sum_c P_c[t],
\label{eq:agg}
\end{equation}
where $O_c$ holds device $c$'s weighted expert outputs and $\mathrm{row}(t,j)$ addresses the output of token $t$'s $j$-th
expert. The arithmetic is the standard local top-$k$ reduction; what \sys contributes is a graph form that a compiler executes
locally and in parallel, independent of the tiers' shapes (Algorithm~\ref{alg:combine}). It makes three choices.

\textbf{Ownership as layout.} A stage or tier computes its $L$ lanes as one output of shape $[L,C,d]$, whose lane axis the
compiler already partitions across the $D$ devices, $L/D$ lanes each. \sys reshapes it, without moving data, to
$[D,(L/D)\,C,d]$, which makes the device the leading batch axis of one batched gather with per-device row indices of shape
$[D,Tk]$: for token $t$'s $j$-th expert $e$, the index on its own device is $l\,C+\mathrm{slot}(e,t)$, with $l$ the expert's
lane on that device, and every other device gets a safe index and a zero mask. Global row indices into the flattened
outputs, or explicit per-device slices, address the same rows, but with them the compiler we use moved expert outputs
between devices (\S\ref{sec:mot-own}).

\textbf{Token-local reduction.} The gathered rows, $[D,Tk,d]$, are summed over $k$ into the partials $[D,T,d]$. The
compiler splits this sum by token tiles, so each core gathers and adds the at most $k$ rows of its own tokens, with no
cross-core reduction. A lane-indexed formulation, which gathers every lane's row for each token and sums over lanes, keeps
device locality but pulls each reduction's operands from all cores (\S\ref{sec:eval-own}).

\textbf{Elementwise cross-device sum.} The partials are summed as $(P_0+P_1)+(P_2+P_3)$ per token tile, which the compiler
lowers as a reduce-scatter with no root device; the next layer's input normalization completes it.

The indices need nothing beyond the routing. The router's top-$k$ output lists each token's experts; integer arithmetic on
the placement and small per-tier tables map an expert to its (tier, device, lane), and the per-expert slot tables that
the routing computes anyway give its row.

\begin{algorithm}[t]
\caption{Ownership-preserving token aggregation for one tier $s$ ($D$ devices)}
\label{alg:combine}
\begin{algorithmic}[1]
\REQUIRE top-$k$ experts $\mathcal{K}[t,1..k]$ of each token $t$; map $\ell(e)\mapsto(s,c,l)$ (tier, device, lane); slot table
  $\mathrm{slot}[e,t]$; tier outputs $O_s$ of shape $[L_s,C_s,d]$, lanes partitioned over devices
\STATE $\hat O_s\leftarrow\mathrm{reshape}(O_s,[D,(L_s/D)\,C_s,d])$ \COMMENT{no data movement}
\FOR{all $t\le T$, $j\le k$, devices $c$ \textbf{in parallel}}
  \STATE $e\leftarrow\mathcal{K}[t,j]$;\; $(s',c',l)\leftarrow\ell(e)$
  \STATE $m[c,t,j]\leftarrow[s'=s\wedge c'=c]$;\; $\mathit{idx}[c,t,j]\leftarrow m\,?\,l\,C_s+\mathrm{slot}[e,t]:0$
\ENDFOR
\STATE $G\leftarrow\mathrm{gather}(\hat O_s,\mathit{idx})\odot m$ \COMMENT{batched over devices: $[D,Tk,d]$}
\STATE $P_s\leftarrow\sum_j G[\cdot,\cdot,j]$ \COMMENT{per token, on its core: $[D,T,d]$}
\STATE \textbf{after all tiers:} $P\leftarrow\sum_s P_s$;\; $y\leftarrow(P_0+P_1)+(P_2+P_3)$ \COMMENT{elementwise}
\end{algorithmic}
\end{algorithm}

\textbf{Cost.} Per device and tier the gather produces $Tk$ rows of $d$ elements, masked slots included, against $L\,T$
rows of a lane-dense accumulator, and a tier adds one such gather whatever its width and capacity; the cross-device traffic
is one $[T,d]$ partial per device. A lane-dense combine instead needs one accumulator per tier width, which grows with $T$
and absorbs the rows the tiers save (Figure~\ref{fig:interaction}).

\subsection{Capacities and Tile-Granular Skipping}
\label{sec:skip}

A compiler may already skip padded work at run time without control flow. Cloud AI 100 lowers the export's mask of empty
slots, $\mathrm{Where}(r<n_e,\,y,\,0)$, into a predicate per row tile of a lane: the tile's gather, GEMMs and SiLU run only
if it holds a routed token, so a lane computes $R\lceil n_e/R\rceil$ rows for tiles of $R$ rows, and an empty expert's
weights are skipped only when its lane is a single tile (Appendix~\ref{app:attrprof}). The tile follows the compiler's
maximum tile size. Such skipping is coarse and does not shrink what a capacity provisions: token tables, buffers and the
weight reads of multi-tile lanes stay at capacity. Rank-bound capacities compose with it: they set the lane size, and with
it the tile, per rank, so the low ranks get single-tile lanes that skip empty experts entirely. Every program therefore
runs at its own best tile size (\S\ref{sec:eval-abl}).
